\section{Experiments}
\label{sec:experiments}

The evaluation asks whether carrying scene state into a new deployment improves
object reconstruction, whether obsolete states are removed, and what happens to
geometric accuracy when earlier coverage is retained. Controlled synthetic data
separates these questions from reference-map incompleteness; public sequences
test additional changes and unchanged scenes. All reported values come from the
same recorded code version. Repeated runs with identical intended behavior vary
by approximately 0.1--0.2 percentage points in the real-room online object score
and 0.1 percentage points in geometry. These are observed run-to-run variations,
not confidence intervals.

\begin{table*}[t]
\centering\small
\caption{Repeated real-room deployments. Online object scores summarize the map
over time. Hidden-change F1 measures changes discovered within a deployment.
Geometry is the room-crop F1 at 5\,cm against that deployment's same-sensor
reference. All values are percentages. Our chained runs import A into B and B
into C; our from-scratch runs import nothing. Dashes are unreported entries, not
zero scores or evidence that a system cannot perform the task.}
\label{tab:real}
\begin{tabular}{llccc}
\toprule
Method & Deployment & Online object P/R/F1 & Hidden-change F1 & Geometry F1 (P/R) \\
\midrule
Ours & A & 100.0/88.7/93.0 & 100.0 & 90.6 (86.0/95.8) \\
Ours & B & 99.4/98.0/98.7 & 100.0 & 92.4 (87.1/98.4) \\
Ours & C & 99.0/98.3/98.6 & 100.0 & 91.6 (85.4/98.9) \\
Ours, from scratch & B & 99.9/93.0/94.6 & -- & 96.6 (95.2/98.0) \\
Ours, from scratch & C & 99.5/95.0/96.4 & -- & 95.2 (92.1/98.4) \\
Khronos & A & 100.0/83.7/89.1 & 100.0 & 91.0 (85.7/97.1) \\
Khronos & B & 96.5/88.8/91.5 & 66.7 & 95.9 (93.5/98.3) \\
Khronos & C & 95.3/77.0/83.8 & 66.7 & 94.7 (91.0/98.7) \\
Panoptic & A & -- & -- & 92.4 (93.8/91.0) \\
Panoptic & B & -- & -- & 93.6 (92.2/95.0) \\
Panoptic & C & -- & -- & 93.0 (88.6/97.8) \\
\bottomrule
\end{tabular}
\end{table*}

\subsection{Data and Comparison Setup}
\label{sec:datasets}

\paragraph{Repeated real-room visits.}
We use three RGB-D deployments, A, B, and C, recorded in one occupied room on
three days. The source dataset is named \emph{room-18}; it contains 20 physical
object identities, changes within every deployment, and 16 annotated changed
states for A-to-B. Each object reference surface is aggregated from depth during
its manually annotated presence interval, using 2\,cm voxels observed in at least
three frames. Scene geometry is evaluated against a separate 1\,cm
reconstruction from that deployment's sensor measurements. This is a
same-sensor reference, not an independent survey of the whole room.
In the reported experimental setup, B and C were reserved for evaluation and
were not used for development.

\paragraph{Controlled synthetic visits.}
Two Isaac Sim sequences revisit a furnished apartment with scripted visible
motion, changes while the camera looks elsewhere, and changes between the end
of A and the start of B. They contain 4,186 and 4,041 frames at 30\,Hz. Object
poses, visibility, and the scene mesh are known from the simulator. The change
annotations contain 24 hidden within-deployment events in A, 10 in B, and 32
cross-deployment events.

\paragraph{Public sequences.}
Three 3RScan reference/rescan pairs~\cite{wald2019rio} provide complementary
cases: an unchanged scene with eight evaluated objects, two moved chairs and a
removed chair, and a scene with a non-rigid blanket change. Each scan is static;
changes occur between scans. Depth is $224\times172$ at 10\,Hz. We initialize
from the reference scan and load its state for the rescan. This is a selected
three-scene evaluation, not a full 3RScan benchmark. Bonn RGB-D
Dynamic~\cite{palazzolo2019bonn} contributes three 15--25\,s box-manipulation
sequences. TUM~\cite{sturm2012tum} supplies walking-person observations; we
report the available quantitative result for \emph{walking\_xyz}.

\paragraph{Compared systems.}
\label{sec:baselines}
We compare the official Khronos release~\cite{schmid2024khronos}, Panoptic
Multi-TSDFs~\cite{schmid2022panoptic}, and our method. Khronos starts each
deployment from an empty map. Panoptic imports its previous map using the same
inter-deployment transform as our method and uses its supplied 3RScan
configuration on that dataset. \emph{Ours, from scratch} runs our own method
without importing any previous-deployment state; it still reconstructs and
updates objects during the current deployment. This isolates the effect of
imported memory, not every individual modification to Khronos. Our backend uses
supplied persistent instance labels; improvements over the unmodified Khronos
pipeline cannot all be attributed to memory alone.

\subsection{What the Scores Measure}
\label{sec:protocols}

\paragraph{Object reconstruction over time.}
The online object score is obtained with a duration-weighted evaluator based on
Khronos, using a 2\,m surface-association gate. Precision (P), recall (R), and
F1 are reported separately as returned by that evaluator. This score measures
object representation over the deployment; it does not test whether an object's
position is correct to 10\,cm. A stale surface 1--2\,m away may still receive
credit. We therefore also report the more local change and residue measurements
below rather than using the online object score alone.

\paragraph{Changes in object state.}
For an annotated object-state surface $G$ and predicted map $M$, let
\begin{equation}
 r(G,M)=\frac{|\{x\in G:d(x,M_{\ell(G)})\le0.10\,\mathrm{m}\}|}{|G|},
 \label{eq:representedness}
\end{equation}
where $M_{\ell(G)}$ contains predicted surface samples of the same semantic
class, $d(x,M_{\ell(G)})$ is the distance to its nearest surface, and $|G|$ is
the number of reference samples. The state is counted as represented when
$r(G,M)\ge0.5$. A moved object
can produce a disappearance of its old state and an appearance of its new state;
event counts are therefore not counts of distinct physical objects.
For within-deployment changes, the test compares states before and after the
annotated change time. For cross-deployment changes, it compares the inherited
scene with the end of the next deployment. Event P/R/F1 are computed from
appeared and disappeared states. \emph{Old-location residue} reports the share
of an annotated departed surface still matched by the final map; smaller values
mean less obsolete geometry remains.

On 3RScan, object identity is matched exactly. Evaluation of a moved object uses
the part of its annotated surface that distinguishes the old and new poses,
rather than their shared surface. Objects require at least 2\,s of visible
frames. The scan and object transforms are composed before assigning motion,
since a listed rigid transformation can correspond to an effectively unchanged
state. On Bonn, the box's reference rest states are constructed from depth,
ground-truth camera poses, and instance masks linked across frames in 3D. These
are derived references, not native object-motion annotations from the dataset.

For Khronos, absence from the rebuilt map can count as disappearance under this
state-comparison protocol even without an explicit deletion decision. This also
produces false disappearances for unchanged objects that a short rescan fails to
reconstruct. Its cross-deployment results must be interpreted with that
asymmetry in mind.

\paragraph{Geometry and contamination.}
Geometry P/R/F1 compares predicted and reference surfaces at 5, 10, and 20\,cm;
the tables report 5\,cm. Real-room comparisons use an oriented crop of the room.
The session-only reference does not contain all previously observed surfaces,
so valid retained geometry can reduce measured precision. We distinguish this
\emph{room-crop score} from an \emph{observed-region score}, which restricts the
comparison to the current deployment's observed domain. The synthetic reference
is the simulator's scene mesh, and 3RScan uses the rescan's reconstruction.
On TUM, contamination is the fraction of final-map vertices that project to
person pixels with agreeing depth in at least three sampled frames. It is a
person-residue proxy, not a full geometric ground truth.

\subsection{Can the Map Retain Coverage and Remove Obsolete States?}
\label{sec:results}

\begin{table}[htbp]
\centering\small
\caption{Cross-deployment changes in the real room. Event scores are percentages;
old-location residue measures departed surface remaining in the final map
(lower is better). Panoptic is run with its previous map loaded.}
\label{tab:real_d3}
\begin{tabular}{llcc}
\toprule
Method & Transition & Event P/R/F1 & Residue \\
\midrule
Ours & A$\rightarrow$B & 100.0/93.8/96.8 & 6.8 \\
Ours & B$\rightarrow$C & 100.0/100.0/100.0 & 10.9 \\
Panoptic & A$\rightarrow$B & 64.0/64.0/64.0 & 57.1 \\
Panoptic & B$\rightarrow$C & 37.5/37.5/37.5 & 60.1 \\
\bottomrule
\end{tabular}
\end{table}

The comparison with our own from-scratch runs in \cref{tab:real} isolates a
benefit of importing scene memory: online object F1 rises from 94.6 to 98.7 in B
and from 96.4 to 98.6 in C. Compared with Khronos, the gains are 7.2 and 14.8
percentage points, but include changes to object-state management as well as
memory. Hidden-change F1 reaches 100.0 in both later deployments, compared with
66.7 for Khronos. The corresponding from-scratch hidden-change scores are not
available in this result set, so this comparison does not isolate memory's
effect on within-deployment change detection.

Earlier coverage need not imply retaining every obsolete state. In
\cref{tab:real_d3}, the method detects 15 of 16 changed states in A-to-B and all
annotated changes in B-to-C. The missed A-to-B case is a swivel chair whose
recorded motion within A and static reference annotation disagree. This case
remains counted as a miss in the reported evaluation. Relative
to Panoptic, old-location residue falls from 57.1\% to 6.8\% and from 60.1\% to
10.9\%. These measurements test actual obsolete surface in the final map,
complementing event recognition.

Geometry reveals a separate tradeoff. Our chained maps achieve recall of
98.4\% and 98.9\%, but room-crop F1 is below both our from-scratch runs and
Panoptic. Two effects contribute: the session-only reference omits some retained
coverage, and imported walls can survive alongside misaligned new copies.
The room-crop scores alone cannot separate these two effects. A complete
per-method observed-region comparison is not available in this result set, so
we do not attribute the entire precision loss to incomplete reference coverage.

\subsection{Controlled Changes with an Exact Scene Reference}

\begin{table}[htbp]
\centering\small
\caption{Synthetic apartment. All scores are percentages. ``Hidden'' denotes
within-deployment hidden-change F1; ``Across'' denotes cross-deployment event F1.
Geometry F1 uses the true scene mesh at 5\,cm.}
\label{tab:synthetic}
\begin{tabular}{llrrrr}
\toprule
Method & Run & Object & Geometry & Hidden & Across \\
\midrule
Ours & A & 84.5 & 97.2 & 100.0 & -- \\
Ours & A$\rightarrow$B & 91.1 & 97.7 & 100.0 & 95.1 \\
Khronos & A & 71.4 & 91.4 & 68.0 & -- \\
Khronos & B & 66.8 & 83.6 & 94.7 & -- \\
Panoptic & A & -- & 93.8 & -- & -- \\
Panoptic & A$\rightarrow$B & -- & 92.1 & -- & 91.8 \\
\bottomrule
\end{tabular}
\end{table}

\begin{table*}[t]
\centering\small
\caption{Selected 3RScan rescans. Scene 1 is unchanged, Scene 2 contains moved
and removed chairs, and Scene 3 contains a non-rigid blanket change. TP/FP/FN
are correctly detected, spurious, and missed change events. Geometry compares
the rescan output to its own reference mesh at 5\,cm. Dashes preserve unreported
or undefined entries in the result records; the counts remain visible.}
\label{tab:3rscan}
\begin{tabular}{llccc}
\toprule
Scene & Method & Event TP/FP/FN & Event P/R/F1 & Geometry F1 (P/R) \\
\midrule
1 & Ours & 0/0/0 & -- & 96.4 (96.8/95.9) \\
1 & Panoptic & 0/1/0 & -- & -- \\
1 & Khronos & 0/3/0 & -- & 94.0 (98.6/89.9) \\
2 & Ours & 4/0/1 & 100/80/88.9 & 90.8 (92.3/89.4) \\
2 & Panoptic & 4/1/1 & 80/80/80.0 & -- \\
2 & Khronos & 4/1/1 & 80/80/80.0 & 90.0 (96.4/84.4) \\
3 & Ours & 0/1/1 & -- & 92.7 (90.0/95.5) \\
3 & Panoptic & 1/2/0 & 33/100/50.0 & -- \\
3 & Khronos & 0/5/1 & -- & 83.4 (98.1/72.6) \\
\bottomrule
\end{tabular}
\end{table*}

With an exact scene reference, our method detects all 24 hidden changes in A and
all 10 in B, as well as 29 of 32 cross-deployment events without a false event
(\cref{tab:synthetic}). Geometry F1 is 97.2 and 97.7, compared with 91.4 and
83.6 for Khronos. Unlike the real-room reference, this comparison is not limited
to surfaces reconstructed by the current sensor trajectory. The missed events
involve a table and sofa whose new locations the camera does not approach, and
a chair displaced 0.6\,m in one frame at the image border. These examples limit
what can be concluded from incomplete observations; they do not show that unseen
changes are reconstructed.

\subsection{Public-Data Checks and Short-Sequence Failures}

\paragraph{Different changes expose different limitations.}
In \cref{tab:3rscan}, the static control produces no false event for our method.
In Scene 2, both moved chairs are represented at their new positions while all
five unchanged objects are retained. The removed chair is missed: the 39\,s
rescan sees through 71\% of its reliable surface, but the accumulated evidence
does not reach the absence threshold. Scene 3 is a failure case, with one missed
and one spurious event; Panoptic detects the annotated change but also reports
two spurious events. Thus these examples do not establish uniform superiority
on all change types. Our geometric F1 is higher than Khronos in all three
rescans, with increased recall and reduced precision. The short scans also
illustrate the cross-deployment scoring asymmetry: rebuilding an unchanged
object incompletely can be scored as a disappearance.

\paragraph{A new deployment must still handle current motion.}
Bonn and TUM do not test cross-deployment memory, but check behavior during one
run. In Bonn's placing sequence, the final box rest state has represented-surface
fractions of 1.00 for our method and 0.43 for Khronos, using
\cref{eq:representedness}. The removing sequence ends one second after pickup,
leaving insufficient evidence for either system. Both retain the first rest
state in the moving sequence: the observed 0.75\,m motion is below Khronos's
1\,m displacement requirement, inherited by our extractor. On TUM
\emph{walking\_xyz}, the person-contamination proxy is 1.02\% for our method
and 1.09\% for Khronos, while our map retains 29\% more static surface than
Khronos. These are useful behavior checks, not a full evaluation of motion
trajectory accuracy.

\FloatBarrier
